\subsection{Local simple-MPDO structure}

% **Normalization boundary:** the raw four-sector tensor in
% thm:mpdo_sal_zcl_rank_one_normalization_obstruction has SAL, literal ZCL, and
% a non-rank-one active trace matrix, but it is not the injective normal tensor
% assumed in Case I. Its normalized representative loses literal ZCL.
% Documented in docs/paper-gaps/cpgsv17_pf_rank_one.tex.
\begin{remark}[Normalization boundary in Lemma C.5]
    \label{rem:mpdo_sal_zcl_rank_one_normalization_boundary}
    Lemma~C.5 of
    \cite[Appendix~C.2, lines~1484--1499]{Cirac2016MPDO_arXiv} claims that an
    injective matrix product density operator tensor satisfying the strong area
    law and zero correlation length has real families $(a_k)_k,(b_k)_k$ with
    \begin{align}
        T_{k,h}=\tr(\eta_{k,h})&=a_kb_h, \notag\\
        \sum_k a_kb_k&=1. \notag
    \end{align}
    The lemma is proved inside Case~I, where the tensor is also assumed to be
    an injective normal tensor. The raw four-sector tensor below has the source
    inverse-map factorization and no such families $a,b$, but it is not the
    normalized tensor of Case~I: its doubled-index transfer satisfies
    $\mathcal E_{\mathcal K}^2=(5/16)\mathcal E_{\mathcal K}$. Its normalized
    representative $A=(4/\sqrt5)\mathcal K$ satisfies
    $\mathcal B_A^2=(4/\sqrt5)\mathcal B_A\ne\mathcal B_A$ and therefore fails
    the paper's literal zero-correlation-length diagram. Thus it does not
    obstruct the all-sector coefficient theorem proved under the complete
    normal Case-I hypotheses. The remaining open step is the structural
    statement used by the following corollary.
\end{remark}

\begin{theorem}[SAL--ZCL rank-one normalization obstruction]
    \label{thm:mpdo_sal_zcl_rank_one_normalization_obstruction}
    \lean{MPOTensor.ActiveSectorSpanningCounterexample.tensor_isSAL}
    \lean{MPOTensor.ActiveSectorSpanningCounterexample.%
      physTraceTransfer_tensor_idempotent}
    \lean{MPOTensor.ActiveSectorSpanningCounterexample.%
      activeSectorTraceMatrix_not_rankOne}
    \lean{MPOTensor.ActiveSectorSpanningCounterexample.%
      tensor_has_sal_sourceZCL_and_non_rankOne_activeTraceMatrix}
    \lean{MPOTensor.ActiveSectorSpanningCounterexample.%
      rectangular_remainder_eq_mul_sub_sq}
    \leanok
    \uses{def:is_sal, def:mpo_source_zcl,
      thm:mpdo_all_cut_markov_implies_sal,
      thm:mpdo_inverse_map_nonzero_rectangular_remainder}
    There is an injective matrix product density operator tensor satisfying
    the strong area law and source zero correlation length whose
    source-selected inverse-map factorization has positive semidefinite
    neighboring operators and a primitive trace-one active trace matrix $T$,
    but for which there are no families $a,b$ satisfying
    $T_{k,h}=a_kb_h$ for all $k,h$. Equivalently, for the corresponding closed
    sector tensors $L,Q$, one has $Q(1-LQ)L=T-T^2\ne0$. Thus the rank-one
    conclusion in Lemma~C.5 fails for this raw representative. The
    normal-tensor hypothesis imposed in Case~I at source lines~1374--1381 is
    not included in this statement.
\end{theorem}

\begin{proof}
    \leanok
    The four physical sectors are one-dimensional. Their cyclic transition
    law is doubly stochastic. After tracing one site, conditioning on the
    middle physical label gives, at every admissible cut,
    \begin{align}
        \rho_{ABC}
        &=\bigoplus_{k=0}^{3}\frac14
          \rho^{L}_{A,k}\otimes\rho^{R}_{k,C}. \notag
    \end{align}
    Thus the reduced state has a four-sector quantum-Markov decomposition and
    the tensor satisfies the strong area law. Its physical-trace
    transfer is the idempotent matrix $LQ$, so it has source zero correlation
    length. The explicit Hayashi decomposition and inverse map recover the
    displayed tensors $L,Q$. The active trace matrix is primitive and has
    trace one, but direct multiplication gives $T^2\ne T$. If
    $T=ab^\top$, then
    \begin{align}
        T^2
        &=a(b^\top a)b^\top
          =\left(\sum_k a_kb_k\right)T
          =T, \notag
    \end{align}
    where the last equality uses $\tr(T)=1$. Hence $T$ has no rank-one
    factorization.
\end{proof}

\begin{theorem}[Two-site blocking preserves renormalization maps]
    \label{thm:mpdo_rfp_via_ts_block_two}
    \lean{MPOTensor.IsRFPViaTS.blockTwo}
    \leanok
    \uses{def:is_rfp_via_ts, def:mpdo_two_site_blocking}
    Suppose that a tensor $K$ has trace-preserving completely positive maps
    $S,T$ satisfying
    \begin{align}
        S[\mathcal K_2(X)]&=\mathcal K_1(X), &
        T[\mathcal K_1(X)]&=\mathcal K_2(X) \notag
    \end{align}
    for every virtual matrix $X$. Then the two-site physical blocking
    $K^{[2]}$ has maps satisfying the same two identities and is therefore a
    renormalization fixed point.
\end{theorem}

\begin{proof}
    \leanok
    Apply $T$ twice to the first physical site, first passing from two sites to
    three and then from three sites to four. This gives a trace-preserving
    completely positive map from the two-site closure to the four-site closure.
    Applying $S$ twice to the first two physical sites in the reverse direction
    gives the coarse-graining map. The canonical identifications of two
    neighboring sites with one blocked site are permutations of orthonormal
    coordinates and preserve both complete positivity and trace preservation.
\end{proof}

\begin{theorem}[Renormalization maps and normalization obstruction for the
    four-sector witness]
    \label{thm:mpdo_four_sector_witness_rfp}
    \lean{MPOTensor.ActiveSectorSpanningCounterexample.normalizedTensor}
    \lean{MPOTensor.ActiveSectorSpanningCounterexample.tensor_isRFPViaTS}
    \lean{MPOTensor.ActiveSectorSpanningCounterexample.%
      blockTwo_tensor_isRFPViaTS}
    \lean{MPOTensor.ActiveSectorSpanningCounterexample.%
      blockTwo_normalizedTensor_not_isRFPViaTS}
    \leanok
    \uses{def:is_rfp_via_ts, def:mpdo_two_site_blocking,
      thm:mpdo_sal_zcl_rank_one_normalization_obstruction,
      thm:mpdo_rfp_via_ts_block_two}
    The four-sector tensor in
    Theorem~\ref{thm:mpdo_sal_zcl_rank_one_normalization_obstruction} has
    trace-preserving completely positive maps $S,T$ such that, for every
    virtual matrix $X$,
    \begin{align}
        S[\mathcal K_2(X)]&=\mathcal K_1(X), &
        T[\mathcal K_1(X)]&=\mathcal K_2(X). \notag
    \end{align}
    Its two-site physical blocking has maps satisfying the same two identities
    and is therefore also a renormalization fixed point. On the other hand,
    for the source-normalized scalar representative
    \begin{align}
        A&=\frac4{\sqrt5}\mathcal K, \notag
    \end{align}
    the two-site blocking $A^{[2]}$ is not a renormalization fixed point.

    The last assertion concerns literal trace-preserving diagrammatic
    renormalization. The rescaled tensor obeys an idempotence relation up to a
    positive scalar, but it does not satisfy the paper's literal
    zero-correlation-length diagram. Consequently it is not, by itself, a
    counterexample to Theorem~4.9 of
    \cite{Cirac2016MPDO_arXiv}. The simple-biCF-BNT hypotheses of that theorem
    are not part of this theorem's statement.
\end{theorem}

\begin{proof}
    \leanok
    Let $\tau=QL$ be the neighboring transition matrix and set
    \begin{align}
        f(i,j)&=(i\bmod 2)+2\left\lfloor\frac j2\right\rfloor. \notag
    \end{align}
    Direct multiplication of the sector matrices gives
    \begin{align}
        S_iS_j&=\tau_{ij}S_{f(i,j)}. \notag
    \end{align}
    For each $k$, the transition weights on the fibre of $f$ satisfy
    \begin{align}
        \sum_{(i,j):f(i,j)=k}\tau_{ij}&=1. \notag
    \end{align}
    The refinement channel has Kraus operators
    $\sqrt{\tau_{ij}}\lvert i,j\rangle\langle f(i,j)\rvert$; the
    coarse-graining channel has Kraus operators
    $\lvert f(i,j)\rangle\langle i,j\rvert$. Their identity resolutions give
    trace preservation, and the two displayed sector identities give the two
    required physical-closure identities. Applying each channel twice and
    regrouping two neighboring physical sites gives the maps for the blocked
    tensor, as in
    Theorem~\ref{thm:mpdo_rfp_via_ts_block_two}.

    Let $P=\operatorname{diag}(1,0)$. The physical-trace transfer of the
    normalized representative is $(4/\sqrt5)P$, while that of its two-site
    blocking is
    \begin{align}
        \mathcal B_{A^{[2]}}&=\frac{16}{5}P. \notag
    \end{align}
    If $A^{[2]}$ admitted the two trace-preserving maps of Definition~4.1,
    trace preservation of the refinement identity would force
    $\mathcal B_{A^{[2]}}^2=\mathcal B_{A^{[2]}}$. This is impossible because
    $P^2=P$ and $16/5\ne1$. The relation
    $\mathcal B_A^2=(4/\sqrt5)\mathcal B_A$ is unchanged in form by a scalar
    normalization when its positive coefficient is allowed to vary. The
    literal diagram in the paper fixes that coefficient to one and is not
    unchanged.
\end{proof}

\begin{theorem}[Inverse-map factorization with a nonzero rectangular remainder]
    \label{thm:mpdo_inverse_map_nonzero_rectangular_remainder}
    \lean{MPOTensor.ActiveSectorSpanningCounterexample.normalizedFourSiteTail_tensor}
    \lean{MPOTensor.ActiveSectorSpanningCounterexample.crossSectorMatrix}
    \lean{MPOTensor.ActiveSectorSpanningCounterexample.sectorMatrix_eq_crossSectorMatrix}
    \lean{MPOTensor.ActiveSectorSpanningCounterexample.crossSectorMatrix_mul}
    \lean{MPOTensor.ActiveSectorSpanningCounterexample.trace_crossSectorMatrix_mul_transfer}
    \lean{MPOTensor.ActiveSectorSpanningCounterexample.threeSiteState}
    \lean{MPOTensor.ActiveSectorSpanningCounterexample.hayashiData}
    \lean{MPOTensor.ActiveSectorSpanningCounterexample.sectorTensorL_hayashiData}
    \lean{MPOTensor.ActiveSectorSpanningCounterexample.sectorTensorR_hayashiData}
    \lean{MPOTensor.ActiveSectorSpanningCounterexample.inverseMapFactorization}
    \lean{MPOTensor.ActiveSectorSpanningCounterexample.%
      inverseMapFactorization_neighboringOperator}
    \lean{MPOTensor.ActiveSectorSpanningCounterexample.rectangular_remainder_ne_zero}
    \lean{MPOTensor.ActiveSectorSpanningCounterexample.%
      inverseMap_provenance_preserves_rectangular_remainder}
    \leanok
    \uses{def:mpdo_zero_weight_reparameterized_inverse_map_factorization,
        thm:mpdo_inverse_map_hayashi_sector_comparison}
    There is a four-sector, bond-dimension-two tensor with a refined Hayashi
    decomposition whose four sectors are one-dimensional and have weight
    $1/4$. Its three-site state is diagonal, with
    \begin{align}
        \rho(i,j,k;i',j',k')
        &=\delta_{(i,j,k),(i',j',k')}
          \frac14(QL)_{i,j}(QL)_{j,k}. \notag
    \end{align}
    The normalized four-site tail is $LQ$. For this tail, the inverse-map
    construction recovers the closed sector tensors $L$ and $Q$ exactly. Its
    neighboring operators therefore coincide with those obtained directly
    from $L$ and $Q$, while $Q(1-LQ)L\ne0$. The inverse-map construction is
    the one in
    \cite[Appendix~C.2, lines~1413--1455]{Cirac2016MPDO_arXiv}.
\end{theorem}

\begin{proof}
    \leanok
    Write $S_i=L_{\mathord\bullet i}Q_{i\mathord\bullet}$ for the four physical
    slices, and let $f_i(a,b)$ be their dual coefficients, so that
    \begin{align}
        \sum_{i=0}^{3}f_i(a,b)S_i&=E_{ab}. \notag
    \end{align}
    With $T=QL$, $p_k=1/4$, and $(LQ)_{00}=1$, the two inverse-map contractions
    are
    \begin{align}
        (LQ)_{00}^{-1}p_k\sum_{i=0}^{3}f_i(0,\beta)T_{ik}
        &=\frac14(4L_{\beta k})=L_{\beta k}, \notag\\
        \sum_{i=0}^{3}f_i(\alpha,0)T_{ki}
        &=Q_{k\alpha}. \notag
    \end{align}
    Hence the neighboring operators are those obtained directly from $L$ and
    $Q$, and direct multiplication gives $Q(1-LQ)L\ne0$.
\end{proof}

% **Normal Case-I boundary:**
% rem:mpdo_sal_zcl_rank_one_normalization_boundary records that the raw witness
% has the displayed SAL, ZCL, and inverse-map identities but is not the normal tensor
% assumed in Case I; its normal representative loses literal ZCL. The completed
% source corollary retains the physical isometry and left--right tensors of the
% selected physical-sector factorization. Documented in
% docs/paper-gaps/cpgsv17_pf_rank_one.tex.
\begin{corollary}[Local sector structure from SAL and ZCL]
    \label{cor:mpdo_sal_zcl_local_sector_structure}
    \lean{MPOTensor.exists_physicalSectorFactorization_rank_one_coefficients_of_isSAL_of_literal_ZCL}
    \leanok
    \uses{def:mpdo, def:injective, def:is_sal,
        def:mpo_physical_trace_transfer, def:nt_cpsv,
        def:mpdo_physical_sector_factorization}
    Let $\mathcal K$ be an injective normal matrix product density operator
    tensor satisfying the strong area law and literal zero correlation length,
    \begin{align}
        \mathcal T_{\mathcal K}^2&=\mathcal T_{\mathcal K}.
        \notag
    \end{align}
    Then there is a physical-sector factorization
    \begin{align}
        \C^d&\simeq\bigoplus_k(B_k^L\otimes B_k^R),
        \notag
    \end{align}
    with
    \begin{align}
        U\kappa_{\beta,\alpha}U^\dagger
        &=\bigoplus_k(l_k)_\beta\otimes(r_k)_\alpha,
        \notag
    \end{align}
    positive semidefinite neighboring operators $\eta_{k,h}$, and real
    families $(a_k)_k,(b_k)_k$ such that
    \begin{align}
        \tr(\eta_{k,h})&=a_kb_h,
        \notag
    \end{align}
    and
    \begin{align}
        \sum_k a_kb_k&=1.
        \notag
    \end{align}
    The orientation of $U$ is the physical-sector-factorization convention
    displayed above. No additional closure property of the sector tensors is
    asserted. This is the structural corollary at
    \cite[Appendix~C.2, lines~1503--1506]{Cirac2016MPDO_arXiv}.
\end{corollary}

\begin{proof}
    \leanok
    \uses{thm:mpdo_casei_all_sector_rank_one_coefficients_of_is_sal}
    Apply Theorem~\ref{thm:mpdo_casei_all_sector_rank_one_coefficients_of_is_sal}.
    Its factorization witness already contains the physical isometry, the
    direct-sum physical-slice formula, and positive neighboring operators.
    Theorem~\ref{thm:mpdo_casei_all_sector_rank_one_coefficients_of_is_sal}
    gives the real-part identity
    $\operatorname{Re}\tr(\eta_{k,h})=a_kb_h$. Positivity makes each
    neighboring operator Hermitian, so its trace is real and the displayed
    complex trace identity follows.
\end{proof}

\begin{theorem}[Positive neighboring sectors generate positive periodic operators]
    \label{thm:mpdo_neighboring_sector_positivity}
    \lean{MPOTensor.PhysicalSectorFactorization.isMPDO_of_neighboringOperator_pos}
    \leanok
    \uses{def:mpdo_physical_sector_factorization}
    Suppose that an MPO tensor $K$ admits a physical-sector factorization
    through an isometry $U$, with neighboring operators
    $\eta_{q,h}$. If $\eta_{q,h}\geq0$ for every pair $(q,h)$, then for every
    $N\geq1$,
    \begin{align}
        \rho^{(N)}(K)&\geq0.
        \notag
    \end{align}
\end{theorem}

\begin{proof}
    \leanok
    In sector coordinates, the periodic operator is block diagonal, and each
    block is a tensor product of positive neighboring operators. Hence every
    block, and therefore the full operator, is positive semidefinite.
    Conjugation by $U^{\otimes N}$ preserves positivity and returns
    $\rho^{(N)}(K)$.
\end{proof}

\begin{definition}[Unstarred orthogonality of BNT layers]
    \label{def:mpdo_bnt_layer_orthogonality}
    \lean{MPOTensor.IsBNTLayerOrthogonal,
        MPOTensor.isBNTLayerOrthogonal_iff_physicalSlice_mul_eq_zero,
        MPOTensor.isBNTLayerOrthogonal_iff_physicalSlice_mul_apply_eq_zero}
    \leanok
    For a family $\{\mathcal K_x\}_{x=0}^{g-1}$, write
    $\mathcal K_x[\beta,\alpha]\in M_d(\C)$ for the physical matrix obtained
    by fixing the two virtual indices. The layers are orthogonal when, for
    $x\ne y$,
    \begin{align}
        \mathcal K_y[\beta_y,\alpha_y]\,
        \mathcal K_x[\beta_x,\alpha_x]&=0
        \label{eq:mpdo_bnt_layer_unstarred}
    \end{align}
    for every choice of virtual indices. Equivalently, for all physical
    indices $i,j$,
    \begin{align}
        \sum_{k=0}^{d-1}
        (\mathcal K_y)^{ik}_{\beta_y\alpha_y}
        (\mathcal K_x)^{kj}_{\beta_x\alpha_x}&=0.
        \notag
    \end{align}
    The factor order is exactly $\mathcal K_y\mathcal K_x$; no adjoint occurs.
    This is equation \texttt{AppKxKy=0} in
    \cite[Appendix~C.2, lines~1634--1639]{Cirac2016MPDO_arXiv}.
\end{definition}

\begin{center}
    \tenkzkernel
    % Exact source orientation: the upper layer is plain K_y and the lower
    % layer is plain K_x.  There is no conjugation or adjoint in AppKxKy=0.
    % Source: Papers/1606.00608/MPDO_KxKy.png and source TeX lines 1634--1639.
    \begin{tenkz}[rows={op,op}, cols=1, west=open, east=open, physical=updown]
        \tn[skin=mpo]{\mathcal K_y} \\
        \tn[skin=mpo]{\mathcal K_x}
    \end{tenkz}
    $ = 0$
\end{center}

\begin{definition}[Physical adjoint]
    \label{def:mpdo_physical_adjoint}
    \lean{MPOTensor.physicalAdjointTensor,
        MPOTensor.physicalSlice_physicalAdjointTensor,
        MPOTensor.mpo_physicalAdjointTensor}
    \leanok
    For an MPO tensor $K$, define $K^\sharp$ by
    \begin{align}
        (K^\sharp)^{ij}_{\beta\alpha}
        &=\overline{K^{ji}_{\beta\alpha}}.
        \notag
    \end{align}
    Thus $K^\sharp[\beta,\alpha]=K[\beta,\alpha]^*$, while the virtual
    indices retain their order. For every periodic chain,
    \begin{align}
        \rho^{(N)}(K^\sharp)_{\sigma,\tau}
        &=\overline{\rho^{(N)}(K)_{\tau,\sigma}}.
        \notag
    \end{align}
\end{definition}

\begin{theorem}[Adjoint closure of the physical slices]
    \label{thm:mpdo_injective_mpdo_adjoint_closure}
    \lean{MPOTensor.IsMPDO.sameMPV₂Pos_physicalAdjointTensor,
        MPOTensor.IsMPDO.sameMPV_physicalAdjointTensor,
        MPOTensor.exists_physicalSlice_conjTranspose_eq_sum_of_isInjective_isMPDO}
    \leanok
    \uses{def:mpdo_physical_adjoint, thm:ft_single}
    Let $K$ be injective and suppose that $\rho^{(N)}(K)\geq0$ for every
    $N\geq1$. Then the linear span of its physical slices is closed under
    conjugate transpose. More precisely, there is $X\in\GL_D(\C)$ such that
    \begin{align}
        K[\beta,\alpha]^*
        &=\sum_{\nu,\mu=0}^{D-1}
          X_{\beta\mu}(X^{-1})_{\nu\alpha}K[\mu,\nu]
        \notag
    \end{align}
    for all virtual indices $\beta,\alpha$.
\end{theorem}

\begin{proof}
    \leanok
    Positivity makes every periodic operator Hermitian, so $K$ and
    $K^\sharp$ generate the same matrix-product vectors. The single-block
    fundamental theorem gives the inner gauge $X$. Taking its virtual matrix
    entries yields the displayed expansion.
\end{proof}

\begin{theorem}[Canonical supports of the active physical-sector factors]
    \label{thm:mpdo_active_physical_sector_factor_supports}
    \lean{MPOTensor.PhysicalSectorFactorization.sectorProductFamily,
        MPOTensor.PhysicalSectorFactorization.IsActiveSector,
        MPOTensor.PhysicalSectorFactorization.isActiveSector_iff_exists_sectorProductFamily_ne_zero,
        MPOTensor.PhysicalSectorFactorization.not_isActiveSector_iff,
        MPOTensor.PhysicalSectorFactorization.not_isActiveSector_iff_left_or_right_eq_zero,
        MPOTensor.PhysicalSectorFactorization.not_isActiveSector_bondDim_zero,
        MPOTensor.PhysicalSectorFactorization.sectorCoordinateTensor_isInjective,
        MPOTensor.PhysicalSectorFactorization.sectorCoordinateTensor_isMPDO_of_neighboringOperator_pos,
        MPOTensor.PhysicalSectorFactorization.sectorProductFamily_conjTranspose_mem_span,
        MPOTensor.PhysicalSectorFactorization.leftTensor_conjTranspose_mem_span_of_isActiveSector,
        MPOTensor.PhysicalSectorFactorization.rightTensor_conjTranspose_mem_span_of_isActiveSector,
        MPOTensor.PhysicalSectorFactorization.leftFactorSupportProj,
        MPOTensor.PhysicalSectorFactorization.rightFactorSupportProj,
        MPOTensor.PhysicalSectorFactorization.activeSector_factorSupport_twoSided,
        MPOTensor.PhysicalSectorFactorization.sectorProductFamily_supportProj,
        MPOTensor.PhysicalSectorFactorization.exists_activeSector_factorSupportIsometries}
    \leanok
    \uses{thm:mpdo_injective_mpdo_adjoint_closure}
    Let $K$ be injective and let
    \begin{align}
        UK[\beta,\alpha]U^*
        =\bigoplus_k l_{k,\beta}\otimes r_{k,\alpha}
        \notag
    \end{align}
    be a physical-sector factorization whose neighboring operators are
    positive semidefinite. Call $k$ active when both families
    $\{l_{k,\beta}\}_\beta$ and $\{r_{k,\alpha}\}_\alpha$ contain a nonzero
    member. For every active $k$, the two factor spans are closed under
    conjugate transpose. If $P_k^L$ and $P_k^R$ are their respective joint
    column-support projections, then
    \begin{align}
        P_k^L l_{k,\beta}
        &=l_{k,\beta}
         =l_{k,\beta}P_k^L,
        &
        P_k^R r_{k,\alpha}
        &=r_{k,\alpha}
         =r_{k,\alpha}P_k^R,
        \notag\\
        \operatorname{supp}
        \{l_{k,\beta}\otimes r_{k,\alpha}\}_{\beta,\alpha}
        &=P_k^L\otimes P_k^R.
        \notag
    \end{align}
    Moreover, there are positive integers $p_k,q_k$ and isometries
    \begin{align}
        V_k^L:\C^{p_k}\longrightarrow\C^{d_k^L},
        \qquad&
        V_k^R:\C^{q_k}\longrightarrow\C^{d_k^R}
        \notag
    \end{align}
    whose range projections are $P_k^L$ and $P_k^R$. If $k$ is inactive,
    every product $l_{k,\beta}\otimes r_{k,\alpha}$ vanishes; in particular,
    there is no active sector when the virtual dimension is zero.
\end{theorem}

\begin{proof}
    \leanok
    The sector-coordinate tensor remains injective. Positivity of the
    neighboring operators makes it a matrix product density operator, so
    Theorem~\ref{thm:mpdo_injective_mpdo_adjoint_closure} closes its physical
    slices under conjugate transpose. Restriction to the $k$-th diagonal
    sector gives coefficients $c_{k,\beta,\alpha}(\mu,\nu)$ such that
    \begin{align}
        (l_{k,\beta}\otimes r_{k,\alpha})^*
        &=
        \sum_{\mu,\nu}c_{k,\beta,\alpha}(\mu,\nu)\,
        l_{k,\mu}\otimes r_{k,\nu}.
        \label{eq:mpdo_active_sector_product_adjoint_closure}
    \end{align}
    Since $k$ is active, choose $\alpha,u,v$ for which
    $r_{k,\alpha}(u,v)\ne0$, and put
    $z=\overline{r_{k,\alpha}(u,v)}$. Evaluating
    (\ref{eq:mpdo_active_sector_product_adjoint_closure}) at the corresponding
    matrix entry and dividing by $z$ yields
    \begin{align}
        l_{k,\beta}^*
        &=
        \sum_\mu
        \left(
            z^{-1}\sum_\nu
            c_{k,\beta,\alpha}(\mu,\nu)r_{k,\nu}(v,u)
        \right)l_{k,\mu}.
    \end{align}
    Interchanging the two factors gives the corresponding identity for
    $r_{k,\alpha}^*$. Thus both factor spans are closed under conjugate
    transpose. The joint column-support projections then absorb the factors
    on both sides. Their tensor product is the support projection of the
    product family, and finite-dimensional range decompositions give the two
    isometries.
\end{proof}

\begin{theorem}[Canonical restriction assembled from the active sector factors]
    \label{thm:mpdo_active_physical_sector_restriction}
    \lean{MPOTensor.PhysicalSectorFactorization.ActiveFactorSupportData,
        MPOTensor.PhysicalSectorFactorization.physicalSupportInclusion,
        MPOTensor.PhysicalSectorFactorization.physicalSupportInclusion_isometry,
        MPOTensor.PhysicalSectorFactorization.%
          physicalSupportInclusion_range_eq_physicalSupportProj,
        MPOTensor.PhysicalSectorFactorization.%
          reindex_physicalSlice_activePhysicalSupportRestriction,
        MPOTensor.PhysicalSectorFactorization.%
          changePhysicalBasis_physicalSupportInclusion_activePhysicalSupportRestriction,
        MPOTensor.PhysicalSectorFactorization.%
          activePhysicalSupportRestrictionData}
    \leanok
    \uses{def:mpdo_physical_support_restriction,
        thm:mpdo_active_physical_sector_factor_supports}
    Let $K$ be injective and let
    \begin{align}
        UK[\beta,\alpha]U^*
        =\bigoplus_k l_{k,\beta}\otimes r_{k,\alpha}
        \notag
    \end{align}
    be a physical-sector factorization whose neighboring operators are
    positive semidefinite.  In every active sector choose isometries onto the
    joint column supports of the left and right factor families, and retain
    no coordinates in an inactive sector.  Their direct sum, transported
    back by $U^*$, is an isometry
    \begin{align}
        V:\C^e\longrightarrow\C^d
        \notag
    \end{align}
    whose range projection is the canonical joint column support
    $P_K$ of the physical slices of $K$.  The restricted tensor
    $K_{\mathrm{act}}=V^*KV$ is injective, satisfies
    \begin{align}
        VK_{\mathrm{act}}V^*&=K,
        \notag
    \end{align}
    and, in the dependent sector coordinates, has physical slices
    \begin{align}
        K_{\mathrm{act}}[\beta,\alpha]
        =\bigoplus_k
        \bigl((V_k^L)^*l_{k,\beta}V_k^L\bigr)
        \otimes
        \bigl((V_k^R)^*r_{k,\alpha}V_k^R\bigr).
        \notag
    \end{align}
    This construction also applies when every sector is inactive: then
    $e=0$, with no complementary summand added.
\end{theorem}

\begin{proof}
    \leanok
    Take the tensor product of the two support isometries in each active
    sector and their dependent direct sum over all sectors, using an empty
    coordinate space in every inactive sector.  Block multiplication gives
    orthonormal columns and the displayed compressed factorization.  The
    range projection absorbs every physical slice.  Conversely, every
    projection that absorbs all physical slices also absorbs the direct sum
    of their joint column supports.  Hence this range projection is $P_K$.
    Two-sided absorption then gives exact re-embedding, and injectivity
    descends through the physical isometry.
\end{proof}

\begin{theorem}[Neighboring positivity on the canonical active restriction]
    \label{thm:mpdo_active_physical_sector_neighboring_positivity}
    \lean{MPOTensor.PhysicalSectorFactorization.FactorActiveSector,
        MPOTensor.PhysicalSectorFactorization.%
          activeRestrictedPhysicalSectorFactorization,
        MPOTensor.PhysicalSectorFactorization.%
          activeRestrictedPhysicalSectorFactorization_neighboringOperator_eq_congruence,
        MPOTensor.PhysicalSectorFactorization.%
          activeRestrictedPhysicalSectorFactorization_neighboringOperator_posSemidef,
        MPOTensor.PhysicalSectorFactorization.ActivePhysicalCompressionWitness,
        MPOTensor.PhysicalSectorFactorization.%
          ActivePhysicalCompressionWitness.factorizationForRestrictionData}
    \leanok
    \uses{thm:mpdo_active_physical_sector_restriction}
    In the setting of
    Theorem~\ref{thm:mpdo_active_physical_sector_restriction}, the restricted
    tensor $K_{\mathrm{act}}$ has a physical-sector factorization indexed
    precisely by the active ambient sectors.  For active sectors $k,h$, put
    \begin{align}
        W_{k,h}=V_k^R\otimes V_h^L.
        \notag
    \end{align}
    Its neighboring operator is
    \begin{align}
        \eta^{\mathrm{act}}_{k,h}
        =W_{k,h}^*\eta_{k,h}W_{k,h}.
        \notag
    \end{align}
    In particular, every $\eta^{\mathrm{act}}_{k,h}$ is positive
    semidefinite.  If the active family is empty, the restricted
    factorization has no sectors and the assertion is vacuous.
\end{theorem}

\begin{proof}
    \leanok
    Remove the empty inactive fibers from the dependent direct sum and
    enumerate the remaining sector subtype by a finite ordinal.  The
    compressed right factor of $k$ and compressed left factor of $h$ give
    \begin{align}
        \sum_\alpha
        \bigl((V_k^R)^*r_{k,\alpha}V_k^R\bigr)
        \otimes
        \bigl((V_h^L)^*l_{h,\alpha}V_h^L\bigr)
        =W_{k,h}^*\eta_{k,h}W_{k,h}.
        \notag
    \end{align}
    Positive semidefiniteness is preserved by this matrix congruence.
\end{proof}

% This theorem supplies the orthogonal-support step used in the completed
% proof of CPSV16 Proposition `prop3to4`. The final proof does not derive or
% assume common-copy-weight independence.
\begin{theorem}[Orthogonal two-sided supports from unstarred layer orthogonality]
    \label{thm:mpdo_bnt_layer_orthogonal_supports}
    \lean{MPOTensor.physicalSliceColumns,
        MPOTensor.physicalSupportProj,
        MPOTensor.physicalSupportProj_mul_physicalSlice,
        MPOTensor.physicalSlice_mul_physicalSupportProj_of_isInjective_isMPDO,
        MPOTensor.physicalSliceColumns_conjTranspose_mul_eq_zero,
        MPOTensor.exists_pairwise_orthogonal_twoSided_physicalSupport,
        MPOTensor.isBNTLayerOrthogonal_of_pairwise_orthogonal_twoSided_support}
    \leanok
    \uses{def:mpdo_bnt_layer_orthogonality,
        thm:mpdo_injective_mpdo_adjoint_closure}
    Let $\{\mathcal K_x\}_{x=0}^{g-1}$ be injective MPO tensors, each of which
    generates positive semidefinite periodic operators at every positive
    length. If, for every pair $x\ne y$, the unstarred identities
    \begin{align}
        \mathcal K_y[\beta_y,\alpha_y]\,
        \mathcal K_x[\beta_x,\alpha_x]&=0
        \notag
    \end{align}
    hold for all virtual indices, then there are pairwise orthogonal
    projections $P_x\in M_d(\C)$ such that
    \begin{align}
        P_x\mathcal K_x[\beta,\alpha]
        &=\mathcal K_x[\beta,\alpha]
         =\mathcal K_x[\beta,\alpha]P_x
        \notag
    \end{align}
    for every $x,\beta,\alpha$. Conversely, any pairwise annihilating family
    with these two-sided support identities satisfies the unstarred layer
    identities.
\end{theorem}

\begin{proof}
    \leanok
    Let $C_x$ be the matrix obtained by stacking the columns of all slices of
    $\mathcal K_x$, and let $P_x$ be the orthogonal projection onto its column
    space. The left support identity is immediate.
    Theorem~\ref{thm:mpdo_injective_mpdo_adjoint_closure} gives the right
    support identity. For $x\ne y$, (\ref{eq:mpdo_bnt_layer_unstarred}) and
    adjoint closure give
    \begin{align}
        C_x^*C_y&=0,
        &P_xP_y&=0.
        \notag
    \end{align}
\end{proof}

\begin{corollary}[Supports of common-weight BNT representatives]
    \label{cor:mpdo_common_weight_bnt_layer_supports}
    \lean{MPOTensor.exists_pairwise_orthogonal_twoSided_physicalSupport_commonWeightAbsorbedBasis}
    \leanok
    \uses{thm:mpdo_neighboring_sector_positivity,
        thm:mpdo_bnt_layer_orthogonal_supports}
    Let a BNT decomposition have a common nonzero weight within each sector,
    and suppose that its one-site words span the product of the virtual matrix
    algebras. For every sector $s$, assume a physical-sector factorization with
    positive neighboring operators. If distinct absorbed representatives obey
    the unstarred layer identities, then there are pairwise orthogonal
    projections $P_s$ such that
    \begin{align}
        P_s\mathcal K_s[\beta,\alpha]
        &=\mathcal K_s[\beta,\alpha]
         =\mathcal K_s[\beta,\alpha]P_s
        \notag
    \end{align}
    for every $s,\beta,\alpha$.
\end{corollary}

\begin{proof}
    \leanok
    One-site spanning makes every absorbed representative injective. The
    positive neighboring operators make every representative an MPDO. Apply
    Theorem~\ref{thm:mpdo_bnt_layer_orthogonal_supports}.
\end{proof}

% **Scope restriction (active physical support compression):** the proof first
% passes through the canonical joint-support restriction, which CPSV16 does
% not form, and then lifts the bond to the printed physical support. See
% `docs/paper-gaps/cpsv16_active_physical_support_compression.tex`.
\begin{theorem}[Positive sector bonds on the printed physical supports]
    \label{thm:mpdo_ambient_sector_factorization_supported_bond}
    \lean{MPOTensor.PhysicalSectorFactorization.%
        exists_etaLocalStructureData_supported_of_twoSidedPhysicalSlice,
        MPOTensor.%
          nonempty_orthogonalCommutingSectorFamily_of_ambientPhysicalSectorFactorization}
    \leanok
    \uses{thm:mpdo_active_physical_sector_neighboring_positivity}
    Let $K$ be injective and admit a physical-sector factorization with
    positive semidefinite neighboring operators.  If an orthogonal projection
    $P$ satisfies
    \begin{align}
        PK[\beta,\alpha]
        &=K[\beta,\alpha]
         =K[\beta,\alpha]P
        \notag
    \end{align}
    for every $\beta,\alpha$, then there is a positive two-site operator $B$
    whose periodic translates commute and such that
    \begin{align}
        (P\otimes P)B(P\otimes P)&=B,
        \notag\\
        \operatorname{mpo}(K,N)
        &=\prod_{j=0}^{N-1}\tau_j(B)
        \qquad (N\geq2).
        \notag
    \end{align}
    Applied to pairwise orthogonal projections $P_s$, this construction gives
    an orthogonal family of positive commuting sector bonds for the tensors
    $K_s$.
\end{theorem}

\begin{proof}
    \leanok
    Restrict $K$ to its canonical joint physical support $Q$ by
    Theorems~\ref{thm:mpdo_active_physical_sector_restriction}
    and~\ref{thm:mpdo_active_physical_sector_neighboring_positivity}.
    The positive neighboring operators on the restriction give a positive
    commuting bond supported on $Q\otimes Q$, with the displayed periodic
    product.  Minimality of the joint column support and the two-sided
    absorption by $P$ give
    \begin{align}
        PQ=Q=QP.
        \notag
    \end{align}
    Hence support of the bond on $Q\otimes Q$ implies support on
    $P\otimes P$.  Performing this construction in every orthogonal sector
    gives the final assertion.
\end{proof}

\begin{theorem}[The five BNT identities give the GSNNCH form]
    \label{thm:mpdo_bnt_physical_sector_gsnnch}
    \lean{MPOTensor.changePhysicalBasis_eq_ite_of_pairwise_orthogonal_twoSided_physicalSupport}
    \lean{MPOTensor.singleKrausMap_sitewise_eq_coefficient_smul_of_orthogonalSupport}
    \lean{MPOTensor.exists_nonnegative_sectorCoefficient_of_orthogonalSupport}
    \lean{MPOTensor.exists_nonnegative_real_eq_sectorCoefficient_of_orthogonalSupport}
    \lean{MPOTensor.%
        hasGSNNCHForm_of_bntLayerOrthogonal_of_physicalSectorFactorization}
    \leanok
    Let $M$ generate positive semidefinite periodic operators and have a BNT
    decomposition with representatives $A_x$ of positive bond dimension,
    copy numbers $n_x\geq1$, and arbitrary copy weights $\mu_{x,q}$.
    Suppose that the decomposition generates the same matrix product vectors
    as $M$ at every positive length.  Assume that the simultaneous one-site
    tuples $(A_x^i)_x$, as $i$ ranges over the physical labels, span the
    product of the virtual matrix algebras.  Assume also that, for distinct
    $x,y$,
    \begin{align}
        A_x[\beta,\alpha]A_y[\beta',\alpha']&=0
        \notag
    \end{align}
    for all virtual indices.  For every $x$, suppose that $A_x$ has the
    physical-sector factorization
    \begin{align}
        U_xA_xU_x^*&=\bigoplus_k l_{x,k}\otimes r_{x,k},
        \notag
    \end{align}
    with positive neighboring operators $\eta^x_{k,h}$ and real numbers
    $a^x_k,b^x_k$ satisfying
    \begin{align}
        \tr(\eta^x_{k,h})&=a^x_kb^x_h,
        &
        \sum_k a^x_kb^x_k&=1.
        \notag
    \end{align}
    Then the matrix product density operators generated by $M$ have the
    GSNNCH form, with the BNT copy numbers $n_x$ as their outer
    multiplicities.  No equality among the copy weights is assumed.
\end{theorem}

\begin{proof}
    \leanok
    \uses{thm:mpdo_ambient_sector_factorization_supported_bond,
        thm:mpdo_bnt_layer_orthogonal_supports,
        thm:mpdo_nonnegative_orthogonal_commuting_sectors_gsnnch}
    The physical-sector identities make every $A_x$ an MPDO.  The unstarred
    layer identities and one-site spanning therefore give pairwise orthogonal
    projections $P_x$ supporting the raw representatives on both sides.
    Fix $N\geq2$ and write
    \begin{align}
        q_x(N)&=\sum_{q=0}^{n_x-1}\mu_{x,q}^{N},
        &
        \rho^{(N)}(M)&=\sum_x q_x(N)\rho^{(N)}(A_x).
        \notag
    \end{align}
    Sitewise compression by $P_x^{\otimes N}$ isolates
    $q_x(N)\rho^{(N)}(A_x)$.  It is positive because $\rho^{(N)}(M)$ is
    positive.  The operator $\rho^{(N)}(A_x)$ is positive and nonzero: the
    latter follows from positive bond dimension and one-site injectivity.
    Hence uniqueness of the positive scalar phase shows that $q_x(N)$ is a
    nonnegative real number.

    Theorem~\ref{thm:mpdo_ambient_sector_factorization_supported_bond} gives
    a supported positive commuting bond $B_x$ with periodic product
    $\rho^{(N)}(A_x)$.  Replace it at this chain length by
    \begin{align}
        \widetilde B_{x,N}
        &=\left(\frac{q_x(N)}{n_x}\right)^{1/N}B_x.
        \notag
    \end{align}
    Then its multiplicity-weighted periodic product is exactly
    $q_x(N)\rho^{(N)}(A_x)$.  The GSNNCH data may depend on $N$, so this
    construction proves the claimed form without copy independence.
\end{proof}

\begin{theorem}[Square-zero obstruction for arbitrary MPO tensors]
    \label{thm:mpdo_bnt_layer_square_zero_obstruction}
    \lean{MPOTensor.BNTLayerOrthogonalityCounterexample.squareZeroPhysicalMatrix_ne_zero,
        MPOTensor.BNTLayerOrthogonalityCounterexample.squareZeroPhysicalMatrix_mul_self,
        MPOTensor.BNTLayerOrthogonalityCounterexample.family_isBNTLayerOrthogonal,
        MPOTensor.BNTLayerOrthogonalityCounterexample.no_pairwise_orthogonal_twoSided_support}
    \leanok
    \uses{def:mpdo_bnt_layer_orthogonality}
    There is a two-label family of bond-dimension-one MPO tensors for which
    the unstarred layer identities hold, but no pairwise orthogonal one-site
    projections support both tensors on both sides. One may take both labels
    to have the single physical slice
    \begin{align}
        A&=|0\rangle\langle1|,
        &A\ne0,
        &A^2=0.
        \notag
    \end{align}
\end{theorem}

\begin{proof}
    \leanok
    The equation $A^2=0$ gives both ordered layer products. If $P_0P_1=0$
    and $P_xA=A$ for both labels, then
    \begin{align}
        A=P_0A=P_0P_1A=0,
        \notag
    \end{align}
    contradicting $A\ne0$.
\end{proof}

If the projections also satisfy $\sum_jP_j=\Id$, then expanding this
identity on both physical indices and using the sector projector relations
leaves only the diagonal term:
\begin{align}
    \mathcal K_i
    &=\sum_{j,j'}P_j\,\mathcal K_i\,P_{j'}
      =P_i\,\mathcal K_i
      =\mathcal K_i\,P_i. \notag
\end{align}
\begin{center}
    \tenkzkernel
    % K_i = sum_{j,j'} P_j K_i P_{j'} = P_i K_i = K_i P_i: each
    % sector tensor absorbs its own projector.  All four terms
    % expose the same open indices -- one west stub, one east
    % stub (K_i's virtual bond), one north stub, one south stub
    % (the physical legs, closed under sum_j P_j = Id) -- so all
    % four panels log the same boundary signature
    % (virtual-west=1 | virtual-east=1 | physical-up=1 | physical-down=1).
    \begin{tenkz}[rows={wire}, cols=1, west=open, east=open, physical=updown]
        \tn[skin=mpo]{K_i}
    \end{tenkz}
    $ = \sum_{j,j'}$
    \begin{tenkz}[rows={op,op,op}, cols=1, west=none, east=none, physical=updown]
        \tn[skin=mpo]{P_j} \\
        \tn[skin=mpo, ports={180:virtual, 0:virtual}]{K_i} \\
        \tn[skin=mpo]{P_{j'}}
    \end{tenkz}
    $ = $
    \begin{tenkz}[rows={op,op}, cols=1, west=none, east=none, physical=updown]
        \tn[skin=mpo]{P_i} \\
        \tn[skin=mpo, ports={180:virtual, 0:virtual}]{K_i}
    \end{tenkz}
    $ = $
    \begin{tenkz}[rows={op,op}, cols=1, west=none, east=none, physical=updown]
        \tn[skin=mpo, ports={180:virtual, 0:virtual}]{K_i} \\
        \tn[skin=mpo]{P_i}
    \end{tenkz}
\end{center}
Equivalently, each local tensor is supported on its own sector:
$\mathcal K_i=P_i\mathcal K_i=\mathcal K_iP_i$.
This is the calculation in
\cite[Appendix~C.2, lines~1745--1751]{Cirac2016MPDO_arXiv}.

% The declaration below isolates the scalar-normalization obstruction. It does
% not assert all standing hypotheses of Theorem 4.9 and therefore does not
% refute that theorem.
\begin{theorem}[Coefficient absorption need not preserve normality]
    \label{thm:mpdo_caseii_absorption_not_normal}
    \lean{MPOTensor.CaseIIAbsorptionCounterexample.printed_absorbed_normality_step_is_false}
    \leanok
    \uses{def:nt_cpsv, def:mpdo_common_sector_weight_absorption,
        def:mpo_physical_trace_transfer}
    There are two bond-one normal tensors $A_0,A_1$ with disjoint doubled
    physical supports and globally normalized weights
    $(\mu_0,\mu_1)=(1/\sqrt2,1)$ such that their weighted bond-diagonal
    assembly has literal physical-trace idempotence. Writing
    $\mathcal E_B$ for the transfer map of a tensor $B$, one has
    \begin{align}
        \rho(\mathcal E_{\mu_0A_0})&=\frac12.
        \notag
    \end{align}
    In particular, the absorbed representative $\mu_0A_0$ is not normal.
\end{theorem}

\begin{proof}
    \leanok
    Let $A_0$ have scalar entries $1/\sqrt2$ on doubled physical symbols
    $(0,0)$ and $(1,1)$, and let $A_1$ have scalar entry $1$ only on
    $(2,2)$.  Their supports are disjoint and, on the one-dimensional virtual
    space,
    \begin{align}
        \mathcal E_{A_0}(x)&=
            \left(\frac12+\frac12\right)x=x,
        & \mathcal E_{A_1}(x)&=x.
        \notag
    \end{align}
    Hence both tensors are normal.  Choose
    $(\mu_0,\mu_1)=(1/\sqrt2,1)$; then $|\mu_s|\leq1$ for both sectors and
    $|\mu_1|=1$, as required by the global normalization.  The weighted
    bond-diagonal assembly has nonzero diagonal physical slices
    \begin{align}
        K^{00}=K^{11}&=\operatorname{diag}(1/2,0),
        & K^{22}&=\operatorname{diag}(0,1),
        \notag
    \end{align}
    so its physical-trace transfer is
    $\mathcal T_K=K^{00}+K^{11}+K^{22}=I$ and therefore
    $\mathcal T_K^2=\mathcal T_K$.  Absorption nevertheless gives
    \begin{align}
        \mathcal E_{\mu_0A_0}
            &=|\mu_0|^2\mathcal E_{A_0}=\frac12\,\mathrm{id},
        & \rho(\mathcal E_{\mu_0A_0})&=\frac12.
        \notag
    \end{align}
    A normal tensor has transfer spectral radius one, so $\mu_0A_0$ cannot be
    normal.
\end{proof}

% CPSV16 Theorem 4.9 (`thm:main-simple`), lines 851--893 and Appendix C.2.
% Provenance: source-label audit completed in issue #5917.
% Gap: docs/paper-gaps/cpgsv17_pf_rank_one.tex.
\begin{theorem}[Printed simple-MPDO implication chain]
    \label{thm:cpsv_theorem49_printed_status}
    \notready
    \uses{def:mpdo_simple_cf_tensor, def:mpdo_two_site_blocking,
        def:is_rfp_via_ts, def:is_sal,
        def:mpdo_source_gsnnch, def:mpo_source_zcl,
        thm:mpdo_sal_zcl_bnt_sector_structure,
        thm:mpdo_sal_zcl_blocking_channels,
        thm:mpdo_caseii_absorption_not_normal}
    Let $K$ be a simple tensor in block-injective canonical form with a basis
    of normal tensors, and let $M=K^{[2]}$.  According to
    \cite[Theorem~4.9]{Cirac2016MPDO_arXiv},
    \begin{align}
        \textnormal{(i)}&\Longrightarrow\textnormal{(ii)}, &
        \textnormal{(ii)}&\Longleftrightarrow\textnormal{(iii)}, &
        \textnormal{(iii)}&\Longrightarrow\textnormal{(iv)}, &
        \textnormal{(iv)}&\Longrightarrow\textnormal{(v)},
        \notag
    \end{align}
    where the conditions are respectively the fixed-point condition, SAL with
    zero correlation length, the commuting Gibbs form with zero correlation
    length, the blockwise local structure, and the fixed-point condition for
    $M$.

    In Case II, the printed implication from SAL and zero correlation length
    to the blockwise local structure applies the normal Case-I factorization to
    each absorbed BNT representative. Absorbing a coefficient $\mu_s$
    multiplies the transfer spectral radius by $|\mu_s|^2$, while the source
    normalization does not imply $|\mu_s|=1$ for every sector.
    Theorem~\ref{thm:mpdo_caseii_absorption_not_normal} shows that this
    absorbed-sector normality route is false under the encoded normalization
    hypotheses. A proof of the implication therefore requires either
    strengthened hypotheses or an alternative factorization argument that does
    not use absorbed-sector normality. This obstruction concerns the printed
    proof route and does not refute Theorem~4.9 itself.
\end{theorem}

% CPSV16 Examples 4.10--4.12, lines 897--942.
% Tracking: issues #5984, #6301, and #6302; Example 4.12 remains under #5919.
% Entropy corrections: docs/paper-gaps/cpsv16_examples_4_10_4_11_entropy.tex.
% Commit a79bfc933 certifies the scalar integer and logarithm inequalities only.
% The spectra, entropy identities, tensor links, and Example 4.12 RFP maps remain open.
\begin{theorem}[Printed mixed-state separation examples]
    \label{thm:cpsv_examples410_412_status}
    \notready
    \uses{def:is_rfp_via_ts, def:is_sal, def:mpdo_source_gsnnch,
        def:mpo_source_zcl, thm:cpsv_example412_literal_formula,
        thm:cpsv_example412_literal_sal,
        thm:cpsv_example412_literal_zcl_rfp_gap}
    The three examples in
    \cite[Examples~4.10--4.12, lines~897--942]{Cirac2016MPDO_arXiv} assert the
    following separations.  Example~4.10 applies correlated on-site bit-flip
    noise to a four-spin ring of Bell pairs and claims zero correlation length
    without saturation of the area law.  The displayed channel repeats the
    left-qubit label; replacing the second occurrence by the right-qubit label
    reproduces all four printed entropies.  At $p=1/4$, the exact calculation
    recorded in the accompanying paper-gap note gives $I_2>I_1$ for the
    displayed reduced spectra.  This calculation alone does not identify those
    density operators as the reductions of the corrected tensor or establish
    its zero-correlation-length property.  The printed universal exception list
    also omits $p=1$, where equality returns.

    Example~4.11 uses the bond-weight matrix
    \begin{align}
        W&=\begin{pmatrix}1&1/2\\1/2&1\end{pmatrix}
        \notag
    \end{align}
    and claims saturation without zero correlation length.  For the literal
    four-site periodic state, the exact calculation recorded in the
    accompanying paper-gap note instead gives $I_2>I_1$ for the displayed
    reduced spectra.  This calculation alone does not identify those density
    operators as reductions of the printed tensor.  Global normalization does
    not change the resulting finite-ring probability distribution.  An open or
    infinite stationary Markov-chain interpretation may explain the source
    claim, but the source does not state that interpretation.

    Example~4.12 uses a two-dimensional tensor generating
    \begin{align}
        \rho^{(N)}&=I^{\otimes N}+\sigma_z^{\otimes N}.
        \notag
    \end{align}
    The source asserts that this tensor has saturation of the area law and zero
    correlation length, is not of the simple commuting form, and is nevertheless
    a renormalization fixed point.  The literal printed tensor is unnormalized:
    its physical-trace transfer squares to twice itself, whereas trace-preserving
    refinement would force literal idempotence.  Thus the fixed-point assertion
    is false for the literal representative.  The tensor $(1/2)M$ is the likely
    intended correction, but its trace-preserving completely positive maps
    remain unconstructed.
\end{theorem}

% CPSV16 Example 4.12, lines 932--939.
% Normalization gap: docs/paper-gaps/cpsv16_example_4_12_normalization.tex.
\begin{definition}[Literal tensor in Example 4.12]
    \label{def:cpsv_example412_literal_tensor}
    \lean{MPOTensor.CPSVExample412Literal.M,
        MPOTensor.CPSVExample412Literal.sigmaZ}
    \leanok
    Let $M$ be the tensor with $d=D=2$ whose only nonzero components are
    \begin{align}
        1=M_{00}^{00}=M_{00}^{11}=M_{11}^{00}=-M_{11}^{11},
        \notag
    \end{align}
    and let
    $\sigma_z=|0\rangle\langle0|-|1\rangle\langle1|$.
\end{definition}

\begin{definition}[One-site Pauli sign]
    \label{def:cpsv_example412_site_sign}
    \lean{MPOTensor.CPSVExample412Literal.siteSign}
    \leanok
    \uses{def:cpsv_example412_literal_tensor}
    For a physical letter $a\in\{0,1\}$, define its Pauli sign by
    \begin{align}
        z_a=(\sigma_z)_{aa},
        \qquad z_0=1,
        \qquad z_1=-1.
        \notag
    \end{align}
\end{definition}

\begin{definition}[Configuration sign]
    \label{def:cpsv_example412_configuration_sign}
    \lean{MPOTensor.CPSVExample412Literal.configurationSign}
    \leanok
    \uses{def:cpsv_example412_site_sign}
    For a binary configuration
    $\boldsymbol\sigma=(\sigma_0,\ldots,\sigma_{N-1})$, define
    \begin{align}
        z(\boldsymbol\sigma)=\prod_{k=0}^{N-1}z_{\sigma_k}.
        \notag
    \end{align}
    This is the diagonal entry of $\sigma_z^{\otimes N}$ at
    $\boldsymbol\sigma$.
\end{definition}

\begin{theorem}[Exact positive-length operator and positivity]
    \label{thm:cpsv_example412_literal_formula}
    \lean{MPOTensor.CPSVExample412Literal.rho_eq_diagonal,
        MPOTensor.CPSVExample412Literal.rho_eq_finKronecker,
        MPOTensor.CPSVExample412Literal.M_isMPDO,
        MPOTensor.CPSVExample412Literal.trace_rho}
    \leanok
    \uses{def:cpsv_example412_literal_tensor,
        def:cpsv_example412_configuration_sign}
    For every $N>0$,
    \begin{align}
        \rho^{(N)}(M)=I^{\otimes N}+\sigma_z^{\otimes N},
        \qquad
        \operatorname{tr}\rho^{(N)}(M)=2^N,
        \notag
    \end{align}
    and $\rho^{(N)}(M)$ is positive semidefinite.
\end{theorem}
\begin{proof}\leanok
    Both tensor powers are diagonal.  On a binary configuration their sum has
    eigenvalue $1+\prod_k z_k$, where each $z_k$ is $1$ or $-1$; hence every
    eigenvalue is either $0$ or $2$.  Summing the signed term over all binary
    configurations gives zero when $N>0$, so the trace is $2^N$.
\end{proof}

\begin{theorem}[One-site loss and saturation]
    \label{thm:cpsv_example412_literal_sal}
    \lean{MPOTensor.CPSVExample412Literal.reducedBlockState_M_eq_scaled_one,
        MPOTensor.CPSVExample412Literal.one_site_trace_loses_sigmaZ,
        MPOTensor.CPSVExample412Literal.M_isSAL}
    \leanok
    \uses{thm:cpsv_example412_literal_formula, def:is_sal}
    Tracing one site removes the $\sigma_z^{\otimes N}$ term.  More generally,
    every nonempty proper normalized $L$-site reduced state is
    $2^{-L}I^{\otimes L}$.
    Therefore its entropy is $L\log 2$, the mutual information is independent
    of $L$, and $M$ satisfies saturation of the area law.
\end{theorem}
\begin{proof}\leanok
    The identity term contributes the dimension of the traced complement,
    while the signed term contains the vanishing one-site sum $1+(-1)$.  The
    entropy of the resulting maximally mixed state is $L\log 2$; the two block
    terms in the mutual information then add to $N\log 2$ independently of the
    cut.
\end{proof}

\begin{theorem}[Scale-two ZCL and the literal RFP obstruction]
    \label{thm:cpsv_example412_literal_zcl_rfp_gap}
    \lean{MPOTensor.CPSVExample412Literal.physTraceTransfer_M,
        MPOTensor.CPSVExample412Literal.physTraceTransfer_M_sq,
        MPOTensor.CPSVExample412Literal.M_isSourceZCL,
        MPOTensor.CPSVExample412Literal.physTraceTransfer_M_not_idempotent,
        MPOTensor.CPSVExample412Literal.M_not_isRFPViaTS}
    \leanok
    \uses{def:cpsv_example412_literal_tensor, def:mpo_source_zcl,
        def:is_rfp_via_ts, thm:rfp_via_ts_physical_trace_idempotent}
    The physical-trace transfer of the printed tensor is
    \begin{align}
        \mathcal T_M=I+\sigma_z
        =\begin{pmatrix}2&0\\0&0\end{pmatrix},
        \qquad
        \mathcal T_M^2=2\mathcal T_M\ne\mathcal T_M.
        \notag
    \end{align}
    Hence the literal representative has source zero correlation length with
    scale $2$, but it cannot satisfy the trace-preserving fixed-representative
    condition.  The tensor $(1/2)M$ is the likely intended correction, but its
    trace-preserving completely positive maps remain unconstructed.
\end{theorem}
\begin{proof}\leanok
    The transfer calculation is immediate from the two diagonal physical
    letters.  Trace preservation of a refinement from one site to two sites
    would force the one-site and two-site traces to agree for every virtual
    boundary, and therefore force $\mathcal T_M^2=\mathcal T_M$, contradicting
    the displayed calculation.
\end{proof}

% **All-sector Case-II boundary:**
% Analytic inheritance is complete for every absorbed BNT representative.
% The explicit bond-one construction formally refutes preservation of
% normality under the source's global weight convention. The final structural
% implication remains not ready; this does not refute Theorem 4.9.
% Documented in
% docs/paper-gaps/cpgsv17_pf_rank_one.tex.
\begin{theorem}[BNT sector structure from SAL and ZCL (all-sector Case-II boundary)]
    \label{thm:mpdo_sal_zcl_bnt_sector_structure}
    \notready
    \uses{def:is_sal, def:mpo_physical_trace_transfer,
        thm:mpdo_sal_bnt_separating_projectors,
        thm:mpdo_absorbed_bnt_sector_caseii_analytic_properties,
        cor:mpdo_sal_zcl_local_sector_structure,
        thm:mpdo_caseii_absorption_not_normal}
    Let $K$ be a simple tensor in block-injective canonical form (biCF),
    equipped with a basis-of-normal-tensors decomposition. Suppose that $K$
    satisfies the strong area law and literal zero correlation length,
    \begin{align}
        \mathcal T_K^2&=\mathcal T_K.
        \notag
    \end{align}
    Then its basis tensors have mutually orthogonal physical supports.
    Moreover, each basis tensor
    has an inverse-map sector decomposition with positive neighboring
    operators $\eta_{k,h}$ and real families $(a_k)_k,(b_k)_k$ satisfying
    \begin{align}
        \tr(\eta_{k,h})&=a_kb_h, \notag\\
        \sum_k a_kb_k&=1. \notag
    \end{align}
    This is implication (ii)$\Rightarrow$(iv) of
    \cite[Theorem~4.9 and Appendix~C.2, lines~1740--1782]
      {Cirac2016MPDO_arXiv}. The separate standing biCF and BNT hypotheses are
    imposed at line~849 and restated at line~1628 of the same source.
    The orthogonal-support conclusion uses the injectivity supplied by biCF,
    the positive neighboring operators in each sector factorization, and the
    unstarred layer identities. The remaining identities specify the local
    factorization and its rank-one normalization.
    Theorem~\ref{thm:mpdo_absorbed_bnt_sector_caseii_analytic_properties}
    proves that every absorbed BNT representative is injective, generates
    positive semidefinite operators, saturates the area law, and retains
    literal physical-trace idempotence. The remaining obstruction is
    normality. If $\mathcal K_s=\mu_sA_s$, then the transfer spectral radius is
    multiplied by $|\mu_s|^2$. The global canonical-form normalization gives
    only one unit-modulus coefficient and does not imply $|\mu_s|=1$ for every
    $s$. Theorem~\ref{thm:mpdo_caseii_absorption_not_normal} gives global
    weights $(1/\sqrt2,1)$, literal ambient physical-trace idempotence, and a
    first absorbed transfer radius equal to $1/2$. Therefore the printed
    assertion that absorption preserves normality is false. The claimed
    all-sector conclusion would instead require stronger assumptions or a
    factorization argument that does not pass through absorbed-sector
    normality. The counterexample isolates this failed intermediate assertion
    and does not contradict the final implication of Theorem~4.9.
\end{theorem}
